% Ad Atticum 7.17 (ad-atticum-07-17) --- English translation
% Translated by Claude (Anthropic) from the Latin (Perseus).
% Style guide: STYLE.md.

\ciceroLetterOpener{CICERO ATTICO salutem}

\ciceroSection{1} Your letter is welcome and a pleasure to me. About transporting the boys to Greece --- I was thinking of that at the time when flight from Italy seemed to be the plan. For we were to make for Spain; this was not equally convenient for them. As for you yourself, along with Sextus, even now you seem to me to be perfectly safe at Rome; for indeed you ought to be the least of friends to our friend Pompey. No one ever knocked so much off the value of city property. Do you see how I am even joking?

\ciceroSection{2} By this time you must know what answers Lucius Caesar is bringing back from Pompey, and what letters from the same man he is taking to Caesar. They were written and given out in such a form as to be displayed in public. On which score I have reproached Pompey in my own mind: a brilliant writer himself, he has handed over the writing of matters of this magnitude --- matters bound to come into everyone's hands --- to our friend Sestius. And so I have never read anything written more \textit{[Greek: in the manner of Sestius]}. All the same, it can be seen from Pompey's letter that nothing is being refused to Caesar, and that everything he asks for is being granted to him in heaped measure. He would be the most insensate of men if he did not accept --- especially given that he made his demands with such utter shamelessness. ``For who are you'' --- so to speak --- ``to dictate that he shall set out for Spain, that he shall disband his garrisons?'' Still, the point is conceded: less honourably now, the commonwealth having been already violated by him and war having been brought against us, than if he had won the right earlier to be a candidate \emph{in absentia}. And even so I am afraid he may not be content with these very terms. For having given those instructions to Lucius Caesar, he ought to have been a little quieter while the answers were being brought back; but it is said that now he is at his most ferocious.

\ciceroSection{3} Trebatius indeed writes that he was asked by Caesar on the ninth day before the Kalends of February to write to me that I should be near the city; that I could do Caesar no greater favour. This at very great length. I gathered from the tally of dates that as soon as Caesar heard of my withdrawal he began to fret that all of us were absent. So I do not doubt that he has written to Piso, and to Servius; what astonishes me is that he has not written to me himself, that he has not approached me through Dolabella, not through Caelius. Still, I do not slight Trebatius's letter --- I know he is uniquely fond of me.

\ciceroSection{4} I wrote back to Trebatius (for to Caesar himself, who had written nothing to me, I was unwilling to write) how difficult that would be at this time; that, however, I was on my own estates and had undertaken no levy and no business. And in this I shall remain so long as there is hope of peace; but if war is waged, I shall not fail my duty or my dignity --- the boys, \textit{[Greek: smuggled out for safekeeping]}, sent on to Greece. For I see that the whole of Italy will be ablaze with war. So much evil has been stirred up, partly by wicked citizens, partly by envious ones. But within a few days it will be plain, from Caesar's answers to our answers, how things are going to come out. Then I shall write you more, if it is to be war; but if --- if instead --- there is even a truce, I shall see yourself, I hope, in person.

\ciceroSection{5} On the fourth day before the Nones of February, the day I gave this letter, I was at my Formian place, where I had returned from Capua, waiting for the women. I had written to them, prompted by your letter, that they should stay at Rome. But I hear there is rather more fear in the city. I wished to be at Capua on the Nones of February, because the consuls had so ordered. Whatever shall be brought here from Pompey, I shall write at once to you; and I shall be expecting your letter on these matters.
